\documentclass[12pt]{article}
\usepackage[utf8]{inputenc}
\usepackage[T1]{fontenc}
\usepackage{hyperref}
\title{Yarsun Language Resources for K-12 Curriculum Development}
\author{Compiled by Claude Sonnet 5.5 (claude-sonnet-5-5) for Andr\'as Kornai}
\date{}
\begin{document}
\maketitle

\section{Introduction}
Yarsun (ISO 639-3 yrs) is listed as an Austronesian (Oceanic, North New Guinea) language of the Sarmi area on Papua's north coast, with about 200 speakers reported in 1999. Its status as a distinct language is uncertain. Wikipedia's article on the Sarmi--Jayapura languages suggests the entry may come from confusing language names with island names. A 2005 SIL note says it may not exist or its speakers may have moved. Searches found no educational materials.

\section{Primary Educational Resources}
None were found.

\section{Linguistic and Reference Materials}
\begin{itemize}
\item Wikipedia, Sarmi--Jayapura languages (discusses the doubtful status of Yarsun): \url{https://en.wikipedia.org/wiki/Sarmi%E2%80%93Jayapura_languages}
\item Endangered Languages Project entry: \url{https://www.endangeredlanguages.com/lang/3166}
\end{itemize}

\section{Religious and Community Resources}
\begin{itemize}
\item Joshua Project profile, Yarsun in Indonesia (a mission-oriented page; its language classification conflicts with the other sources): \url{https://joshuaproject.net/people_groups/15996}
\end{itemize}

\section{Digital and Audio Resources}
None were found.

\section{Curriculum Development Context}
Before any curriculum work, fieldwork would be needed to confirm whether Yarsun exists as a separate language, where it is spoken and by whom. Nothing in the search results supports planning school materials yet.
\end{document}